\section{Physical windows from a separating cut}\label{s:cut}

We use the shells of Definition~\ref{d:shells}, the plane
specifications of Definition~\ref{d:framemodel}, and the original-leg
partition sums of Definition~\ref{d:pinmodel}.  In a sum with subscript
\(G\setminus Y\), the original volume remains \(G\): the retained
legs come from \(\mathcal L(G)\), and the protected layer is always
\(S(G)\), with no new pins at removed sites.  The only response
assumption in this section is \(\mathsf P(A,\theta)\) of
Definition~\ref{h:P}.  Throughout, \(d\) is ambient \(\ell^1\)
distance.  Distances in the square geometry, including the distances
to corners in the shell definition, are ambient \(\ell^\infty\)
distances.  An exterior half-plane is understood through its lattice
sites; for example, the left exterior half-plane of \([-u,u]^2\) is
\(\{z\in\mathbb Z^2:z_1\le-u-1\}\).

\subsection{Shell geometry}

The parallel-side assertion below uses the same nonexceptional index
condition as the bound on \(A_i\).  This condition is necessary.  For
example, take \(l_0=64\), an integer \(s\ge390\), \(u=s^4\), \(a=0\),
\(y=(2s-u,-u-1)\), and \(i=2\).  The left side of the radius-\(k_2\)
square has first coordinate \(-u\), meets \(C_u(0)\), and has distance
one from the left exterior half-plane.  The corner \((-u,-u)\)
belongs to \(D_2\), so this index is exceptional.

\begin{lemma}\label{l:shellgeom}
For \(u\ge M_{\mathrm g}(l_0)=\max\{16,6(l_0+1)\}^4\),
\(|I_0|\le1\), and \(D_i\subset F_i\) for \(2\le i\le i_1\).
For \(2\le i\le i_1\) outside \(I_0\), \(|A_i|\le8l_0\).
For the same indices \(2\le i\le i_1\) outside \(I_0\), every side of
\(y+[-k_i,k_i]^2\) parallel to an edge of \(C_u(a)\)
and meeting \(C_u(a)\) is at distance at least \(s/3\) from the
exterior half-plane beyond that edge.
For \(i\ge2\), every point of \(C_u(a)\setminus F_i\)
within \(\ell^\infty\) distance one of \(F_i\) lies in \(D_{i-1}\);
\(y\) is not within that distance of \(F_i\).
\end{lemma}

\begin{proof}
Translate the coordinates so that \(a=0\).  We may interchange the
coordinates and reverse their signs for this geometric proof.  Since
\(\operatorname{dist}_\infty(y,C_u(0))=1\), we can consequently write
\[
 y=(\beta,-u-1),\qquad -u-1\le\beta\le u+1.
\]
These changes of coordinates concern sets and distances only.
They make no assertion about the frame specification.

We have \(u=s^4\), \(s\ge16\), and \(s\ge6(l_0+1)\).  In particular,
\[
 s^2+\tfrac43s<u,\qquad s^2+s<u/2,\qquad
 \lfloor s\rfloor>2s/3,\qquad s/3>l_0.
\]
For example, \(s^2+s\le(17/16)s^2<s^4/2\),
\(s^2+4s/3\le(13/12)s^2<s^4\), and
\(\lfloor s\rfloor\ge s-1>2s/3\).
For \(2\le i\le i_1\), we also have \(k_i\le s^2+s\).

Both upper corners of \(C_u(0)\) have distance \(2u+1\) from
\(y\).  At least one lower corner has distance at least \(u\)
from \(y\), because
\( |\beta+u|+|\beta-u|\ge2u\).
Thus at most one corner can have distance strictly less than \(u\)
from \(y\).  If a corner \(c\) has distance less than \(s/3\)
from a point \(z\in D_i\), then
\[
 |c-y|_\infty\le |c-z|_\infty+k_i
       <s^2+\tfrac43s<u.
\]
Only this one possible corner can therefore cause an exceptional
index.  For distinct indices the radii differ by at least
\(\lfloor s\rfloor>2s/3\).  Two points lying within \(s/3\)
of the same corner have distance strictly less than \(2s/3\),
whereas points of the two different spheres have distance at least
the difference of their radii.  This proves \(|I_0|\le1\).

Fix now \(2\le i\le i_1\), \(i\notin I_0\), and put \(k=k_i\).
The bottom horizontal side of the radius-\(k\) square misses
\(C_u(0)\).  Its top side has second coordinate \(-u-1+k\).
If this side meets \(C_u(0)\), its distances from the bottom and
top exterior half-planes are respectively \(k\) and \(2u+2-k\).
They are at least \(s\) and greater than \(u\), respectively,
since \(k\ge2s-1\ge s\) and \(k<u/2\).

Suppose the left vertical side meets \(C_u(0)\).  Its first
coordinate is \(\beta-k\in[-u,u]\), and its lowest point in
\(C_u(0)\) is \((\beta-k,-u)\in D_i\).  Its distance to the
lower left corner is \(\beta-k+u\).  Nonexceptionality gives
\(\beta-k+u\ge s/3\), so the whole side has distance
\(\beta-k+u+1\ge s/3+1\) from the left exterior half-plane.
Its distance from the right exterior half-plane is
\(u+1-\beta+k\ge k\ge s\).
For the right vertical side the lowest point is
\((\beta+k,-u)\); its distance to the lower right corner is
\(u-\beta-k\).  Thus its distance from the right exterior
half-plane is at least \(s/3+1\), and its distance from the
left exterior half-plane is \(\beta+k+u+1\ge k\ge s\).
This proves the parallel-side assertion, including when
\(\beta=\pm(u+1)\).

A point of \(A_i\) has distance at most \(l_0\) from at least
one of the four exterior half-planes.  A shell side parallel to
that half-plane cannot contain this point, by the clearance just
proved.  On a perpendicular side at most \(l_0\) lattice sites
have that distance at most \(l_0\).  There are four shell sides
and two perpendicular half-planes for each side.  Counting with
possible repetitions therefore gives \(|A_i|\le8l_0\).

For every \(2\le i\le i_1\),
\(k_i\ge is-1>(i-1)s\), proving \(D_i\subset F_i\).
Finally let \(i\ge2\), \(x\in C_u(0)\setminus F_i\), and
\(z\in F_i\) with \(|x-z|_\infty\le1\).  With \(r=(i-1)s\),
integrality and the Lipschitz property of \(|\,\cdot-y|_\infty\)
give
\[
 \lfloor r\rfloor\le |z-y|_\infty-1
       \le |x-y|_\infty\le\lfloor r\rfloor.
\]
Consequently \(x\in D_{i-1}\).  Each point of \(F_i\) has
distance greater than \((i-1)s\ge s>1\) from \(y\).
If \(F_i\) is empty both final assertions are vacuous, with
distance to it equal to \(+\infty\).
\end{proof}

\subsection{Factorization over a cut}

\begin{lemma}\label{l:cut}
Let \(F\subset\mathbb Z^2\) be finite, \(B\subset F\) have no
neighbour outside \(F\), and \(G=F\setminus B\).
Suppose \(X\subset G\) and \(G\setminus X=G_1\sqcup G_2\), with
no bond between the pieces, no site of \(G_1\) adjacent to \(F^c\),
and no site of \(G_2\) adjacent to \(B\).
For the original legs of \(G\), use pins \(\eta\) into \(F^c\)
and window letters \(s\) into \(B\). Put
\[
 e(\eta,s)=\log\Psi_{G\setminus X}(\eta s)-\log\Psi_G(\eta s).
\]
If \(|e(\eta,s)|\le E<\infty\) for all admissible \(\eta,s\),
then there is a strictly positive probability \(\nu\) on window
configurations, independent of \(\eta\), such that
\[
 e^{-2E}\nu(s)\le\mu^{+,\eta}_{\kappa,F}(n_B=s)
                         \le e^{2E}\nu(s).
\]
Consequently any two window laws have total variation at most
\(1-e^{-2E}\).
One may take \(\nu(s)\) proportional to \(W_B(s)Z_{G_1}^{s}\),
where \(W_B\) is the actual product of frame site weights and
internal window bond weights. Only original legs are retained in
the partition functions after removing \(X\).
\end{lemma}

\begin{proof}
Fix a proper colouring \(\kappa\).  The window configuration space
is the nonempty finite product \(\Omega_B=\prod_{x\in B}D_x\).
For \(s\in\Omega_B\) define
\[
 W_B(s)=\prod_{x\in B}\frac1{12}
       \prod_{\substack{\{x,z\}\subset B\\d(x,z)=1}}
                       (1+s_x\cdot s_z).
\]
All these weights are positive by Lemma~\ref{l:threeframes}.
The absence of bonds from \(B\) to \(F^c\) implies, on summing
the actual frame weights over \(G\), that
\begin{equation}\label{e:windowfactor}
 \mu^{+,\eta}_{\kappa,F}(n_B=s)
       =\frac{W_B(s)Z_G^{\eta s}}
                    {Z^{+,\eta}_{\kappa,F}}.
\end{equation}
Here the sum over \(G\) equals \(Z_G^{\eta s}\) by
Lemma~\ref{l:exact}: every summed spin has at most four incident
bond or original-leg factors.  This use of quadrature concerns
only a partition polynomial, with the letters \(s\) held fixed.
All partition sums in this proof are finite and strictly positive
by Lemma~\ref{l:finitefacts}; the actual frame normalizer is also
strictly positive by Definition~\ref{d:framemodel}.

For \(r=1,2\), define the partition sums on \(G_r\) by retaining
only the original legs of \(G\) whose tails lie in \(G_r\).
In particular, do not introduce factors on legs into \(X\).
The hypotheses on adjacency say that the retained pins on \(G_1\)
are all window letters and those on \(G_2\) are all exterior
letters.  The absence of bonds between the pieces and the product
site functional give
\begin{equation}\label{e:cutfactorization}
 Z_{G\setminus X}^{\eta s}=Z_{G_1}^{s}Z_{G_2}^{\eta},
 \qquad
 Z_{G\setminus X}^{0}=Z_{G_1}^{0}Z_{G_2}^{0}.
\end{equation}
The empty-piece convention makes these identities valid when
one or both pieces are empty.  Using the definition of \(e\),
we obtain
\[
 Z_G^{\eta s}
   =e^{-e(\eta,s)}
       \frac{Z_G^0 Z_{G_2}^{\eta}}{Z_{G_1}^0Z_{G_2}^0}
       Z_{G_1}^{s}.
\]
Set
\[
 Q=\sum_{s\in\Omega_B}W_B(s)Z_{G_1}^{s},\qquad
 \nu(s)=Q^{-1}W_B(s)Z_{G_1}^{s}.
\]
Every term in \(Q\) is positive and \(\Omega_B\) is nonempty,
so \(0<Q<\infty\), and \(\nu\) is strictly positive and does
not depend on \(\eta\).  Equations~\eqref{e:windowfactor}
and~\eqref{e:cutfactorization} imply
\[
 \mu^{+,\eta}_{\kappa,F}(n_B=s)
       =c(\eta)\nu(s)e^{-e(\eta,s)},\qquad
 c(\eta)^{-1}=\sum_{s\in\Omega_B}\nu(s)e^{-e(\eta,s)}.
\]
Admissible boundary and window assignments exist, so the assumed
absolute-value bound also implies \(E\ge0\).  The last sum lies
in \([e^{-E},e^E]\); its reciprocal lies in the same interval.
This proves the two pointwise comparisons.

For any two finite probability laws \(p,p'\),
\[
 d_{\mathrm{TV}}(p,p')
    =1-\sum_s\min\{p(s),p'(s)\}.
\]
Both window laws are at least \(e^{-2E}\nu\), so their overlap
sum is at least \(e^{-2E}\).  The asserted total variation bound
follows.  When \(B=\varnothing\), \(\Omega_B\) has its single
empty configuration, \(\nu\) assigns it mass one, and the same
proof gives total variation zero.
\end{proof}

\subsection{The cost of removing a connected set}

\begin{lemma}\label{l:removal}
Assume \(\mathsf P(A,\theta)\), and put
\(\mathsf F(d)=A\theta^{\max\{d,0\}}\).
Let \(x_1,\ldots,x_p\) be distinct points of \(G\), \(p\ge1\),
with \(b_G(x_j)=d(x_j,G^c)\ge3\), such that each \(x_j\), \(j\ge2\),
has a nearest neighbour among its predecessors.
For \(X=\{x_1,\ldots,x_p\}\) and every original unit-pin assignment,
\[
 |\log\Psi_{G\setminus X}(\tau)-\log\Psi_G(\tau)|
 \le1+4\mathsf F(b_G(x_1)-3)
                  +2\sum_{j=2}^p\mathsf F(b_G(x_j)-3).
\]
\end{lemma}

\begin{proof}
Here \(G\) is finite, as in Definition~\ref{d:pinmodel}.
Put \(Y_0=\varnothing\), \(Y_j=\{x_1,\ldots,x_j\}\), and
\(D_j=G\setminus Y_j\).  The depth assumption implies
\(Y_j\subset G\setminus S(G)\).  At the \(j\)-th removal let
\(u_1,\ldots,u_k\) be the surviving neighbours of \(x_j\).
The site \(x_j\) has no original leg, all four of its original
neighbours belong to \(G\), and \(k=4\) for \(j=1\).
For \(j\ge2\) one predecessor is a neighbour, so \(0\le k\le3\).

Integrate the spin at \(x_j\) in the finite functional.  By
Lemma~\ref{l:moments}, the resulting polynomial is
\begin{equation}\label{e:removalpolynomial}
 \begin{split}
 I_j&=\begin{cases}
 1,&k=0,1,\\
 1+\frac13\displaystyle\sum_{1\le r<t\le k}n_{u_r}\cdot n_{u_t},
       &k=2,3,\\
 1+\frac13\displaystyle\sum_{1\le r<t\le4}n_{u_r}\cdot n_{u_t}
                    +\frac1{15}Q_j,&k=4,
 \end{cases}\\
 Q_j&=(n_{u_1}\cdot n_{u_2})(n_{u_3}\cdot n_{u_4})
      +(n_{u_1}\cdot n_{u_3})(n_{u_2}\cdot n_{u_4})
      +(n_{u_1}\cdot n_{u_4})(n_{u_2}\cdot n_{u_3}).
 \end{split}
\end{equation}
The displayed definition of \(Q_j\) is made only in the case
\(k=4\); no \(u_4\) or quartic polynomial is formed in the other
cases.  For either
\(\bullet=\tau\) or \(\bullet=0\), finite summation gives
\(Z_{D_{j-1}}^\bullet=Z_{D_j}^\bullet
\mathbb E_{D_j}^\bullet I_j\).  Positivity of all normalizers
follows from Lemma~\ref{l:finitefacts}.  Consequently
\begin{equation}\label{e:removaltelescope}
 \log\Psi_G(\tau)-\log\Psi_{G\setminus X}(\tau)
   =\sum_{j=1}^p
      \left(\log\mathbb E_{D_j}^\tau I_j
                         -\log\mathbb E_{D_j}^0 I_j\right).
\end{equation}

We first give a uniform lower bound on each expectation in
\eqref{e:removaltelescope}.  Distinct neighbours of a square-lattice
site have distance two from one another, so no two \(u_r\)'s are
adjacent.  Each satisfies \(b_G(u_r)\ge b_G(x_j)-1\ge2\),
has no original leg, and has at most three present bonds after
\(x_j\) is removed.  Condition on every spin except these
\(k\) spins, on a conditioning event of positive probability.
Because there is no bond between these sites and none has an original
leg, the joint conditional weight is a constant times the product,
over \(r\), of its site weight and its surviving bond factors.
Summing over the \(k\) spins gives the product of their one-site
normalizers, each positive by Lemma~\ref{l:finitefacts}.
The joint conditional law is therefore the product of these
one-site laws.  Lemma~\ref{l:35} applies to each of them, with at
most three factors, so their means \(m_r\) satisfy \(|m_r|\le3/5\).
For \(k=0\) the product is the empty product.  Thus
\[
 \sum_{r<t}m_r\cdot m_t
    =\frac12\left(\left|\sum_r m_r\right|^2
                                  -\sum_r|m_r|^2\right)
    \ge-\frac{k}{2}\frac9{25}.
\]
For \(k\le3\) this proves
\(\mathbb E_{D_j}^\bullet I_j\ge1-3k/50\ge41/50>1/2\).
For \(k=4\), the pointwise bound \(|Q_j|\le3\) gives
\[
 \mathbb E_{D_j}^\bullet I_j
       \ge1-\frac6{25}-\frac3{15}
       =\frac{14}{25}>\frac12.
\]
For \(k=0\) the polynomial is one and its expectation is one.
These arguments average over positive conditioning fibres only;
zero-probability fibres are immaterial.

For each pair \(r<t\), the sites \(u_r,u_t\) survive in
\(D_j\), have distance two, and have the common removed neighbour
\(x_j\in Y_j\).  If \(z\in S(G)\), the distance function
\(b_G\) is one-Lipschitz and \(b_G(z)\le2\).  Therefore
\[
 d(u_r,z)\ge b_G(u_r)-2\ge b_G(x_j)-3,
 \qquad d(u_r,S(G))\ge b_G(x_j)-3.
\]
Property \(\mathsf P(A,\theta)\) now gives
\[
 \left|\mathbb E_{D_j}^\tau(n_{u_r}\cdot n_{u_t})
             -\mathbb E_{D_j}^0(n_{u_r}\cdot n_{u_t})\right|
       \le \mathsf F(b_G(x_j)-3).
\]
The set \(S(G)\) here is the original one throughout the
telescoping.  In particular, no distance to a newly created hole
has entered this estimate.

For positive real \(a,b\ge1/2\), integration of \(1/t\)
between them gives \(|\log a-\log b|\le2|a-b|\).
When \(j\ge2\), there are at most three pairs in
\eqref{e:removalpolynomial}, and hence the absolute value of
the \(j\)-th logarithmic difference is at most
\(2\mathsf F(b_G(x_j)-3)\).  At \(j=1\) there are six pairs,
and \(|\mathbb E_{D_1}^\tau Q_1-\mathbb E_{D_1}^0Q_1|\le6\).
The corresponding logarithmic difference is at most
\[
 2\left(2\mathsf F(b_G(x_1)-3)+\frac6{15}\right)
       =4\mathsf F(b_G(x_1)-3)+\frac45
       \le4\mathsf F(b_G(x_1)-3)+1.
\]
All inserted pair and four-spin products used here have separate
degree at most four: their endpoints have lost the bond to \(x_j\)
and each endpoint occurs once in an insertion.  No higher moment
identity is used.  Summing the bounds in
\eqref{e:removaltelescope} proves the claim, including \(p=1\).
\end{proof}

\subsection{A tapered set of physical sites}\label{s:cutgeom}

\begin{definition}\label{d:thickcut}
Fix \(l_0=64\), a nonexceptional shell with \(2\le i\le s\),
and write \(k=k_i\), \(F=F_i\), \(B=B_i\), \(G=F\setminus B\).
For \(z\in\mathbb Z^2\), set
\[
 t(z)=\big||z-y|_\infty-k\big|,\qquad
 h(z)=\operatorname{dist}_\infty(z,C_u(a)^c)\quad
        (h(z)=0\text{ outside }C_u(a)).
\]
Define
\[
 \begin{gathered}
 U=\{z\in C_u(a):t(z)\le\min(h(z)-32,s/12)\},\\
 X_0=\{z\notin U:\operatorname{dist}_\infty(z,U)=1\},\qquad
 X=\{z:d(z,X_0)\le1\}.
 \end{gathered}
\]
Set \(G_1=(U\setminus B)\setminus X\) and
\(G_2=G\setminus(U\cup X)\).
\end{definition}

\begin{lemma}\label{l:cutgeom}
There exists \(s_0\ge\max\{16,6(64+1)\}\) such that, uniformly
in the translate, \(y\), and nonexceptional integer \(2\le i\le s\),
the following hold for \(s=u^{1/4}\ge s_0\).
The set \(U\) is finite, nonempty, and star-connected, and \(U^c\)
is star-connected. We have \(B\subset U\subset F\), all neighbours
of \(B\) lie in \(U\), and no site of \(U\) is adjacent to \(F^c\).
The set \(X_0\) is nonempty and lies in \(G\). For \(x\in X_0\),
\[
 b_G(x)\ge
 \begin{cases}\max\{t(x),16\}-2,&t(x)\le s/12,\\s/12-1,&t(x)>s/12,
 \end{cases}
 \qquad b_G(x)\ge14 .
\]
For each integer \(j\ge0\), there are at most \(48\) sites of
\(X_0\) with \(t(x)=j\le s/12\).
For \(0<\vartheta<1\), put
\[
 C_{\rm near}(\vartheta)=48\vartheta^{14}
                \left(17+\frac1{1-\vartheta}\right).
\]
Then
\[
 \sum_{x\in X_0}\vartheta^{b_G(x)}
 \le C_{\rm near}(\vartheta)
                 +4(2u+1)^2\vartheta^{s/12-1}.
\]
The thickened set \(X\) lies in \(G\), has \(b_G\ge13\), and
\(G\setminus X=G_1\sqcup G_2\) satisfies all separation and
no-adjacency assumptions of Lemma~\ref{l:cut}. Moreover
\[
 \sum_{z\in X}\vartheta^{b_G(z)}
       \le5\vartheta^{-1}\sum_{x\in X_0}\vartheta^{b_G(x)} .
\]
\end{lemma}

\begin{proof}
We may take \(s_0=390\).  Translate so that \(a=0\), and put
\(C=C_u(0)\), \(T=s/12\), and \(\rho(z)=|z-y|_\infty\).
Here \(s,T\) are real, whereas \(k,\rho(z),t(z),h(z)\) are
integers at lattice sites.
The functions \(h(z)=\max\{0,u+1-|z|_\infty\}\), \(\rho\),
and \(t=|\rho-k|\) are one-Lipschitz for \(\ell^\infty\)
distance.  We shall repeatedly use the following consequences of
\(s\ge390\), \(u=s^4\), and \(2\le i\le s\):
\begin{equation}\label{e:sitecutscales}
 \begin{gathered}
 2s-1\le k\le s^2\le u-31,\qquad T\ge1,\qquad
 k-T-1>(i-1)s,\\
 s/3-2T=s/6\ge65>32,\qquad
 2u+2-2T>32,\qquad u+1-T>32.
 \end{gathered}
\end{equation}
Indeed, the third inequality follows from
\(k\ge is-1\) and \(s-T-2=11s/12-2>0\).
The inequalities involving \(u\) follow from \(u=s^4\)
and \(s\ge390\); for instance \(s^4-s^2>31\) and
\(s^4-s/12>32\).  All constants in what follows are thus uniform.

\emph{Connectivity of the tapered set.}
On all of \(\mathbb Z^2\) its definition can be written as
\begin{equation}\label{e:taperconditions}
 U=\{z:t(z)\le T,\ h(z)-t(z)\ge32\}.
\end{equation}
The second condition forces \(z\in C\).
If \(\rho(z)>k\), move each coordinate whose absolute
displacement from \(y\) is maximal one step toward the corresponding
coordinate of \(y\).  This is a star-step, decreases \(\rho\)
and \(t\) by one, and decreases \(h\) by at most one.
If \(\rho(z)<k\), move one coordinate with maximal absolute
displacement one step away from \(y\); this is a nearest-neighbour
step, increases \(\rho\) by one, and again decreases \(t\) by
one and \(h\) by at most one.  The case of zero displacement can
use either sign; in fact \(y\notin U\).  Starting in \(U\),
these steps preserve both inequalities in
\eqref{e:taperconditions}.  After exactly \(t(z)\) steps they
reach
\[
 \mathcal K=\{z:\rho(z)=k,\ h(z)\ge32\}\subset U.
\]

We verify that \(\mathcal K\) is nonempty and nearest-neighbour
connected.  Set \(m=u-31\).  Since \(|y_j|\le u+1=m+32\)
and \(32\le k\le m\), each integer interval
\[
 J_j=[y_j-k,y_j+k]\cap[-m,m]\cap\mathbb Z,
       \qquad j=1,2,
\]
is nonempty.  It contains at least one of \(y_j-k,y_j+k\):
otherwise \([-m,m]\cap\mathbb Z\) would lie strictly between
these two integers, which is impossible since \(m\ge k\).
Consequently the rectangle \(J_1\times J_2\) contains a corner
of \(y+[-k,k]^2\).  Reverse the signs of the coordinates for
this geometric argument so that the contained corner is
\(y-(k,k)\).  The intersections of this rectangle with the left
and bottom sides of the square are nonempty lattice segments
joined at that corner.  If the right side meets the rectangle,
its bottom endpoint \(y+(k,-k)\) lies on the bottom segment.
If the top side meets the rectangle, its left endpoint
\(y+(-k,k)\) lies on the left segment.  Thus all side segments
are joined by nearest-neighbour paths.  Their union is precisely
\(\mathcal K\), since \(h\ge32\) is equivalent to
\(|z|_\infty\le m\).  The radial paths above now show that
\(U\) is nonempty and star-connected.  It is finite because
it is contained in \(C\).

The exterior \(C^c\) is nearest-neighbour connected.  To see this
directly, a site with first coordinate outside \([-u,u]\) can
move vertically to height \(u+1\), remaining outside \(C\).
A site whose second coordinate is outside \([-u,u]\) can first
move horizontally to first coordinate \(u+1\), then vertically
to height \(u+1\).  Along the row at height \(u+1\) these paths
join.  We join each \(z\in C\setminus U\) to \(C^c\) while
avoiding \(U\).  There are three cases.
If \(\rho(z)>k+T\), repeatedly move a coordinate of maximal
displacement away from \(y\).  Each step increases \(\rho\),
so stays outside \(U\), and the path eventually leaves \(C\).
If \(\rho(z)<k-T\), repeatedly move a coordinate of maximal
nonzero displacement toward \(y\).  The value of \(\rho\)
does not increase, and the sum of the absolute displacements
decreases by one at every step.  The path stays outside \(U\)
and reaches \(y\notin C\), unless it leaves \(C\) earlier.
In the remaining case \(t(z)\le T\) and \(h(z)-t(z)<32\).
Here \(z\ne0\), since \(h(0)-t(0)\ge u+1-T>32\) whenever
\(t(0)\le T\).  Move a coordinate of maximal absolute value
away from zero and continue in that direction until leaving
\(C\).  At every step in \(C\), \(h\) decreases by one,
whereas \(t\) decreases by at most one.  Thus \(h-t<32\)
persists, independently of whether \(t\le T\) persists.
This path also avoids \(U\).  These cases prove that \(U^c\)
is nearest-neighbour connected, hence star-connected.

\emph{The window and the raw boundary.}
For \(z\in B\), \(t(z)=0\) and \(h(z)\ge65\), so \(z\in U\).
A nearest neighbour of such a point has \(t\le1\) and
\(h\ge64\); it too belongs to \(U\).  For \(z\in U\) we
have \(h(z)\ge32\) and \(\rho(z)\ge k-T\).  Hence \(z\in F\),
and every nearest neighbour \(w\) of \(z\) satisfies
\(h(w)\ge31>0\) and \(\rho(w)\ge k-T-1>(i-1)s\).
There is therefore no bond from \(U\) to \(F^c\).

If \(x\in X_0\), choose \(v\in U\) with
\(|x-v|_\infty=1\).  Then
\(h(x)\ge31>0\) and \(\rho(x)\ge k-T-1>(i-1)s\),
so \(x\in F\).  Since \(x\notin U\) and \(B\subset U\),
we have \(x\in G\).  The raw boundary is finite, and is
nonempty: a point of \(U\) with largest first coordinate has
its right nearest neighbour in \(X_0\).

\emph{Depths of the raw boundary.}
Fix \(x\in X_0\), put \(t=t(x)\), \(h=h(x)\), and keep the
point \(v\) chosen above.  Distance from \(x\) to \(B\) in
the \(\ell^\infty\) metric is at least \(t\), and also at
least \(65-h\), by the Lipschitz properties of \(\rho\) and
\(h\).  Its distance to \(C^c\) is \(h\).  For
\(w\in C\setminus F\),
\[
 |x-w|_\infty\ge\rho(x)-\rho(w)
       \ge k-t-(i-1)s\ge s-1-t.
\]
If this last set is empty the same lower bound is valid with
distance \(+\infty\).  As \(G^c=B\cup C^c\cup(C\setminus F)\)
and \(d\ge\operatorname{dist}_\infty\), we conclude that
\begin{equation}\label{e:sitecutdepth}
 b_G(x)\ge
    \min\{\max(t,65-h),\ h,\ s-1-t\}.
\end{equation}

Suppose first \(t\le T\).  Since \(x\notin U\) but \(v\in U\),
\[
 h<32+t,\qquad
 h\ge h(v)-1\ge32+t(v)-1\ge30+t.
\]
Thus
\begin{equation}\label{e:rawdepthrange}
                30+t\le h<32+t.
\end{equation}
In \eqref{e:sitecutdepth}, the first entry is at least
\(\max(t,33-t)\), which is at least \(\max(t,16)\).
The second entry is at least \(30+t\ge\max(t,16)-2\).
The third is at least \(11s/12-1\), which is at least
\(\max(t,16)-2\): use \(t\le T\), \(s\ge390\), and
\(s-1-2T>0\).  This proves the stated near depth.

If \(t>T\), then \(t\le T+1\), and
\(h\ge32+t(v)-1\ge30+t>T+30\).
The three entries in \eqref{e:sitecutdepth} are consequently
at least \(T\), \(T+30\), and \(11s/12-2\), respectively.
Each is at least \(T-1\), proving the far depth.
The near depth is at least \(14\), and the far depth is
at least \(s/12-1\ge390/12-1>14\).

\emph{Counting near sites, including possible coincident sides.}
Fix an integer \(0\le j\le T\).  A raw boundary site with
\(t(x)=j\) lies on a square boundary of radius \(k-j\) or
\(k+j\) about \(y\), and on a boundary \(h=h'\), where
\(h'\) is an integer in \([30+j,32+j)\).
There are at most three choices for \(h'\).  These levels
are the boundaries of squares about zero of radius \(u+1-h'\),
which is positive by \eqref{e:sitecutscales}.

We exclude coincident supporting lines of these two square
boundaries.  For each edge of \(C\), define the signed inward
depth of a parallel line.  For a vertical line with first
coordinate \(q\), the left and right depths are
\(q+u+1\) and \(u+1-q\), respectively; for horizontal lines
use the same expressions in the second coordinate.  Each side
of a radius-\(k\pm j\) square corresponds to the side with
the same orientation and sign of the radius-\(k\) square.
Their signed depths from any fixed parallel exterior half-plane
differ by at most \(j\).

The varying-coordinate intervals of both corresponding sides
meet \([-u,u]\): their centres have coordinate in
\([-u-1,u+1]\), and their radii are at least
\(k-j\ge2s-1-T\ge1\).
If the original side's fixed coordinate lies in \([-u,u]\),
the original side therefore meets \(C\).  We have
\(u\ge390^4=M_{\mathrm g}(64)\), \(2\le i\le s<i_1\), and
\(i\notin I_0\).  Thus the nonexceptional parallel-side assertion
of Lemma~\ref{l:shellgeom} applies and gives original signed depth
at least \(s/3\).  The shifted depth is at least
\[
                   s/3-j>32+j>h'.
\]
The strict inequality uses \(s/3-2j\ge s/6>32\).
If the original line lies beyond the specified edge, its signed
depth is at most zero; the shifted depth is at most
\(j<30+j\le h'\).  If it lies beyond the opposite edge,
its original depth is at least \(2u+2\); the shifted depth is
at least \(2u+2-j>32+j>h'\), by
\eqref{e:sitecutscales}.  These three cases exhaust all positions
of the original line.  None can yield shifted depth \(h'\),
which would be necessary for coincidence with a side of the
\(h'\)-square.

Without coincident supporting lines, a side of one axis-parallel
square meets the boundary of another in at most two points:
its parallel sides do not intersect it, and there are only two
perpendicular sides.  Their square boundaries thus meet in
at most eight lattice sites.  There are at most two radius
choices and at most three depth levels.  We have proved the
required bound \(2\cdot3\cdot8=48\).

For \(0<\vartheta<1\), the near sites therefore contribute at most
\[
 \begin{split}
 48\sum_{j=0}^{\infty}\vartheta^{\max(j,16)-2}
 &=48\vartheta^{14}
                    \left(17+\frac{\vartheta}{1-\vartheta}\right)\\
 &\le48\vartheta^{14}
                    \left(17+\frac1{1-\vartheta}\right).
 \end{split}
\]
The far sites contribute at most
\(|X_0|\vartheta^{T-1}\le(2u+1)^2\vartheta^{T-1}\),
which is at most the displayed far term in the statement.
This proves the raw cut sum.

\emph{Thickening and separation.}
The ambient distance function \(b_G(z)=d(z,G^c)\), defined
also outside \(G\) by zero, is one-Lipschitz for \(d\).
Every \(z\in X\) has a point \(x\in X_0\) with \(d(z,x)\le1\).
It follows that \(b_G(z)\ge b_G(x)-1\ge13\), so \(X\subset G\).
In particular \(X\) meets neither \(B\) nor the protected
original layer \(S(G)\).  Directly summing over the five sites
in each nearest-neighbour cross
gives
\[
 \sum_{z\in X}\vartheta^{b_G(z)}
    \le\sum_{x\in X_0}\sum_{d(z,x)\le1}
                              \vartheta^{b_G(z)}
    \le5\vartheta^{-1}\sum_{x\in X_0}\vartheta^{b_G(x)}.
\]

Since \(B\subset U\subset F\) and \(X\subset G\), the two
sets specified in Definition~\ref{d:thickcut} give the disjoint
union \(G\setminus X=G_1\sqcup G_2\).  If a bond joined
\(p\in G_1\) to \(q\in G_2\), then \(p\in U\), \(q\notin U\),
and \(|p-q|_\infty=1\); hence \(q\in X_0\subset X\), a
contradiction.  A point of \(G_1\subset U\) is not adjacent to
\(F^c\), as already proved.  A point adjacent to \(B\) lies
in \(U\), so cannot belong to \(G_2\).  Finally \(B\) itself
has no neighbour in \(F^c\), because \(B\subset U\) and
\(U\) has this property.  These are all the separation and
no-adjacency assumptions of Lemma~\ref{l:cut}.
\end{proof}

\begin{lemma}\label{l:Xconn}
Under Lemma~\ref{l:cutgeom}, the finite nonempty cut \(X\) is
nearest-neighbour connected. Hence it can be ordered so that every
site after the first has a previously listed nearest neighbour.
\end{lemma}

\begin{proof}
Lemma~\ref{l:cutgeom} says that \(U\) is finite, nonempty and
star-connected, and that \(U^c\) is star-connected.
Lemma~\ref{l:boundaryconn} in Appendix~\ref{s:topology} therefore makes
\(X_0=\partial_\infty U\) nonempty and star-connected.
Its finiteness follows from that of \(U\).
Lemma~\ref{l:thicken} in that appendix makes \(X\) nearest-neighbour connected,
and its definition makes it finite and nonempty.

For completeness, choose a root \(x_1\in X\) and order all
points by nondecreasing shortest-path distance from the root
in the nearest-neighbour graph induced by \(X\), breaking
ties arbitrarily.  These distances are finite nonnegative
integers.  Every nonroot point has a neighbour at distance
one smaller, obtained from a shortest path; that neighbour
has already been listed.  This also covers a singleton cut.
\end{proof}

\subsection{The uniform window estimate}

\begin{theorem}\label{t:cutMOS}
If \(\mathsf P(A,\theta)\) holds, then the family of plane frame
specifications has \(\mathsf W(64,p_0,M_\star)\), with common
\(p_0\in(0,1]\) and integer \(M_\star\ge1\) depending only on
\(A,\theta\). The constants work for every proper plane
three-colouring and every translate.
Specifically, let
\[
 E=1+4A+10A\theta^{-4}
       \left[48\theta^{14}\left(17+\frac1{1-\theta}\right)+1\right],
 \qquad p_0=e^{-2E}.
\]
There exists an integer \(M_\star\ge1\) for which (W) holds
with this displayed choice of \(p_0\).
\end{theorem}

\begin{proof}
Fix \(A\ge2\) and \(0<\theta<1\) satisfying
Definition~\ref{h:P}.  We first choose a common scale threshold.
Let \(c=-\log\theta>0\).  For \(u=s^4\ge1\),
\[
 4(2u+1)^2\theta^{s/12-1}
       \le36\theta^{-1}s^8 e^{-cs/12}
       \le K_\theta s^{-2},\qquad
 K_\theta=36\theta^{-1}\,10!\left(\frac{12}{c}\right)^{10}.
\]
The second inequality follows from
\(e^r\ge r^{10}/10!\) for \(r=cs/12>0\), by the
nonnegative power series for the exponential.
Choose, for example,
\[
 M_\star=\left\lceil\max\{390^4,K_\theta^2,1\}\right\rceil.
\]
Then every integer \(u\ge M_\star\) satisfies the threshold
of Lemma~\ref{l:cutgeom} and
\begin{equation}\label{e:windowtail}
                 4(2u+1)^2\theta^{s/12-1}\le1.
\end{equation}

Fix any proper plane colouring, any translate \(C_u(a)\),
any indicated exterior site \(y\), and any nonexceptional
integer \(2\le i\le s\).  Form \(F,B,G,U,X_0,X,G_1,G_2\)
as in Definition~\ref{d:thickcut}.  Lemmas~\ref{l:cutgeom}
and~\ref{l:Xconn} give a nonempty finite cut in \(G\), an
ordering \(x_1,\ldots,x_p\) as in Lemma~\ref{l:removal},
and \(b_G\ge13\) on it.  For every admissible exterior
assignment \(\eta\) and window assignment \(s_B\), use the
unit pins \(\eta s_B\) on the original legs of \(G\).
The removal lemma applies to these pins without any restriction
on their directions.  As \(b_G-3\ge0\), it yields
\begin{equation}\label{e:windowcutcost}
 \begin{split}
 |e(\eta,s_B)|
 &\le1+4A\theta^{b_G(x_1)-3}
                    +2A\theta^{-3}\sum_{j=2}^p\theta^{b_G(x_j)}\\
 &\le1+4A+2A\theta^{-3}\sum_{z\in X}\theta^{b_G(z)}\\
 &\le1+4A+10A\theta^{-4}
       \left[48\theta^{14}\left(17+\frac1{1-\theta}\right)
                        +4(2u+1)^2\theta^{s/12-1}\right]
 \le E.
 \end{split}
\end{equation}
The first site is kept as a seed in this estimate; its additive
cost is included in the term one.

All assumptions of Lemma~\ref{l:cut} hold by
Lemma~\ref{l:cutgeom}.  That lemma supplies one strictly
positive reference law on the actual window alphabet and the
minorization \(e^{-2E}\nu\), uniformly in \(\eta\).
Thus any two admissible boundaries for \(F_i\) give window
marginals with total variation at most \(1-e^{-2E}\).
The statement is also automatic for an empty window, whose
marginal is a point mass.  This is precisely
\(\mathsf W(64,p_0,M_\star)\) in Definition~\ref{d:W}.
All terms defining \(E\) are finite and positive, so
\(p_0=e^{-2E}\in(0,1]\).

The geometric estimates did not use a colouring, and the
removal estimate was uniform over original unit pins.
Quadrature in Lemma~\ref{l:cut} transfers only the sums over
the sites outside the window.  The reference law retains
\(W_B\) with the chosen frames and may vary with the colouring
and placement.  No equality of window laws at different
colourings or translates is needed.  The constants chosen
above consequently apply to every colouring and translate.
\end{proof}
